\subsection{Adaptive Gradient Methods}

\textbf{First-order adaptive methods} like Adam~\citep{kingma2015adam}, RMSprop~\citep{tieleman2012rmsprop}, and AdaGrad~\citep{duchi2011adagrad} adapt learning rates per parameter using gradient statistics. AdamW~\citep{loshchilov2019adamw} improved Adam by decoupling weight decay, while AMSGrad~\citep{reddi2018amsgrad} addressed convergence issues via maximum variance tracking. However, all share fundamental limitations: (1) conflating noise with curvature in second moments, (2) narrow learning rate ranges, (3) warmup requirements.

\textbf{AdaBelief}~\citep{zhuang2020adabelief} attempted to measure ``prediction error'' via $(\g_t - \m_t)^2$, but this quantity is corrupted: it equals $\beta_1(\g_t - \m_{t-1})$ at steady state---a scaled, biased version of innovation. Critically, AdaBelief uses this term \emph{inversely} (in denominator), causing division-by-zero instabilities and severe overfitting in large-batch settings (our CIFAR-100 experiments: -2\% accuracy vs IRIS despite lower training loss).

\subsection{Momentum and Extrapolation Methods}

\textbf{Nesterov momentum}~\citep{nesterov1983method} uses gradient at extrapolated position. \textbf{Adan}~\citep{xie2022adan} extends this with gradient difference $d_t = \g_t - \g_{t-1}$ for adaptive Nesterov momentum. However, gradient difference has fundamental flaws:
\begin{itemize}
    \item \textbf{High variance:} $\Var[d_t] = \Var[\g_t] + \Var[\g_{t-1}] \approx 2\sigma^2$ (two noisy samples)
    \item \textbf{Vanishing correction:} Near minima, $\g_t \approx \g_{t-1} \implies d_t \approx 0$ (no stabilization)
    \item \textbf{Convergence-divergence cycles:} Our experiments show catastrophic mid-training spikes (Rosenbrock: $10^{-4} \to 10^{-1}$ at step 400; Rastrigin: repeated $10^6\times$ spikes)
\end{itemize}

IRIS addresses these issues by using innovation $\Innovation_{t|t-1} = \g_t - \ghat_{t-1}$ (gradient vs smooth estimate), reducing variance to $\sigma^2$ and maintaining active correction even when gradients vanish.

\subsection{Second-Order Methods}

\textbf{Newton-type methods} like L-BFGS~\citep{liu1989lbfgs}, K-FAC~\citep{martens2015kfac}, and Shampoo~\citep{gupta2018shampoo} use curvature information but incur prohibitive computational costs ($O(d^2)$ or worse). \textbf{AdaHessian}~\citep{yao2021adahessian} approximates diagonal Hessian efficiently but still requires Hessian-vector products. \textbf{Sophia}~\citep{liu2023sophia} clips by Hessian diagonal but exhibits aggressive clipping (70\% engagement) and 2$\times$ per-step cost.

IRIS achieves second-order-like adaptivity at first-order cost: innovation variance $(\g_t + \beta_2 \Innovation_{t|t-1})^2$ implicitly captures curvature without explicit Hessian computation.

\subsection{Kalman Filtering in Optimization}

Extended Kalman filters~\citep{kalman1960filter} provide optimal state estimation under Gaussian noise. Recent work on Kalman-inspired optimizers~\citep{vuckovic2018kalman,olivier2021kalman} focused primarily on state estimation rather than error correction. IRIS's key innovation is using the \emph{innovation sequence} (prediction errors) as a first-class signal for gradient correction and variance estimation, inspired by Kalman filtering's innovation analysis but adapted to stochastic optimization's unique challenges.
